\section{Pendientes}

\subsection{Sin autorización externa}

\begin{itemize}[leftmargin=*,itemsep=0.2em]
\item Construir y probar la imagen Docker y el arranque completo.
\item Validación final del formulario contra el papel del colegio; los pasos
  2--5 se comprobaron por HTTP el 23-11-2026.
\item Ejecutar los 14 casos pendientes del plan manual: 11 de la hoja 04
  (medicación TENS) y TRX-11/12/13.
\item Cargar alumnos reales (script \codigo{limpiar-para-carga-manual.sql}
  listo, sin ejecutar).
\item Manual de uso con capturas reales (la maqueta está en
  \path{documentacion/manual-usuarios/}).
\end{itemize}

\subsection{Despliegue en el servidor}

Aplicar pasos \codigo{12}, \codigo{13}, migración \codigo{AgregarSso}, guarda
de privilegios, fuentes PDF y release actual. Respaldos programados y copia
externa, rotación de credenciales visibles, mover la base a
\codigo{/srv/mpaz-data} y publicación con dominio propio y Cloudflare Access.

\subsection{Decisiones del colegio y del equipo}

Textos legales de autorizaciones, matriz final de permisos, protocolo de
medicación TENS (MAT-05), paracetamol, formato del certificado de alumno
regular, conservación de datos, revisión jurídica y Ley 21.719; formato
ministerial y quién lo sube, identificador común con Asistencia (MAT-15),
enlace desde la página del colegio, consumo del SSO por otros módulos y
alcance del registro público.

\nota{Este documento se revisó contra el código el 27-11-2026: lo descrito
es lo que la aplicación hace hoy. Los ítems pendientes no se documentan como
funcionalidad existente.}
